\section{V-trace}
\label{sec:v-trace}

Off-policy learning is important in the decoupled distributed actor-learner architecture because of the lag between when actions are generated by the actors and when the learner estimates the gradient. To this end, we introduce a novel off-policy actor-critic algorithm for the learner, called V-trace. 

First, let us introduce some notations. We consider the problem of discounted infinite-horizon RL in Markov Decision Processes (MDP), see \cite{Puterman:1994,sutton-barto98} where the goal is to find a policy $\pi$ that maximises the expected sum of future discounted rewards:
$V^{\pi}(x) \eqdef\E_{\pi}\big[\sum_{t\geq 0} \gamma^t r_t\big]$, where $\gamma\in[0,1)$ is the discount factor, $r_t = r(x_t, a_t)$ is the reward at time $t$, $x_t$ is the state at time $t$ (initialised in $x_0=x$) and $a_t\sim\pi(\cdot|x_t)$ is the action generated by following some policy $\pi$. 

The goal of an off-policy RL algorithm is to use trajectories generated by some policy $\mu$, called the {\em behaviour policy}, to learn the value function $V^{\pi}$ of another policy $\pi$ (possibly different from $\mu$), called the {\em target policy}. 

\subsection{V-trace target}
Consider a trajectory $\left(x_t, a_t, r_t\right)_{t=s}^{t=s+n}$ generated by the actor following some policy $\mu$. We define the $n$-steps V-trace target for $V(x_s)$, our value approximation at state $x_s$, as:
\beqa
v_s &\eqdef& \textstyle{V(x_s) + \sum_{t=s}^{s+n-1} \gamma^{t-s} \Big( \prod_{i=s}^{t-1} c_i \Big ) \delta_t V}, \label{eq:target.off}
\eeqa
where $\delta_t V \eqdef \rho_t \big(r_t+\gamma V(x_{t+1})-V(x_t)\big)$  is a temporal difference for $V$, and  $\rho_t\eqdef\min\big(\bar\rho, \frac{\pi(a_t|x_t)}{\mu(a_t|x_t)}\big)$ and $c_i\eqdef\min\big(\bar c, \frac{\pi(a_i|x_i)}{\mu(a_i|x_i)}\big)$ are truncated importance sampling (IS) weights (we make use of the notation $\prod_{i=s}^{t-1} c_i=1$ for $s=t$). In addition we assume that the truncation levels are such that $\bar \rho\geq \bar c$.

Notice that in the on-policy case (when $\pi=\mu$), and assuming that $\bar c \geq 1$, then all $c_i=1$ and $\rho_t=1$, thus   (\ref{eq:target.off}) rewrites
\begin{align}
v_s &= V(x_s) + \textstyle{\sum_{t=s}^{s+n-1}} \gamma^{t-s} \big(r_t+\gamma V(x_{t+1})-V(x_t)\big) \notag \\
&= 
\label{eq:target.on}
\textstyle{\sum_{t=s}^{s+n-1}} \gamma^{t-s} r_t + \gamma^{n} V(x_{s+n}),
\end{align}
which is the on-policy $n$-steps Bellman target. Thus in the on-policy case, V-trace reduces to the on-policy $n$-steps Bellman update. This property (which Retrace \cite{munos2016safe} does not have) allows one to use the same algorithm for off- and on-policy data. 
 
Notice that the (truncated) IS weights $c_i$ and $\rho_t$ play different roles.  The weight $\rho_t$ appears in the definition of the temporal difference $\delta_t V$ and defines the fixed point of this update rule. In a tabular case, where functions can be perfectly represented, the fixed point of this update (i.e., when $V(x_s)=v_s$ for all states), characterised by $\delta_t V$ being equal to zero in expectation (under $\mu$), is the value function $V^{\pi_{\bar \rho}}$ of some policy $\pi_{\bar \rho}$, defined by
 \beq\label{eq:pi_rho}
 \pi_{\bar \rho}(a|x)\eqdef \frac{\min \big(\bar \rho \mu(a|x),\pi(a|x)\big)}{\sum_{b\in A}\min \big(\bar \rho \mu(b|x),\pi(b|x)\big)},
 \eeq
  (see the analysis in Appendix~\ref{sec:analysis}
).
So when $\bar \rho$ is infinite (i.e.~no truncation of $\rho_t$), then this is the value function $V^{\pi}$ of the target policy.  However if we choose a truncation level $\bar \rho<\infty$, our fixed point is the value function $V^{\pi_{\bar \rho}}$ of a policy $\pi_{\bar\rho}$ which is somewhere between $\mu$ and $\pi$. At the limit when $\bar \rho$ is close to zero, we obtain the value function of the behaviour policy $V^{\mu}$. In Appendix~\ref{sec:analysis}
we prove the contraction of a related V-trace operator and the convergence of the corresponding online V-trace algorithm.

The weights $c_i$ are similar to the ``trace cutting" coefficients in Retrace. Their product $c_s\dots c_{t-1}$ measures how much a temporal difference $\delta_tV$ observed at time $t$ impacts the update of the value function at a previous time $s$. The more dissimilar $\pi$ and $\mu$ are (the more off-policy we are), the larger the variance of this product. We use the truncation level $\bar c$ as a variance reduction technique. However notice that this truncation does not impact the solution to which we converge (which is characterised by $\bar \rho$ only).

Thus we see that the truncation levels $\bar c$ and $\bar \rho$ represent different features of the algorithm:  $\bar \rho$ impacts the nature of the value function we converge to, whereas $\bar c$ impacts the speed at which we converge to this function.
 
\begin{remark}
V-trace targets can be computed recursively:
$$v_s = V(x_s) + \delta_sV + \gamma c_s\big( v_{s+1} - V(x_{s+1})\big).$$
\end{remark}
\begin{remark}
Like in Retrace($\lambda$), we can also consider an additional discounting parameter $\lambda\in [0,1]$ in the definition of V-trace by setting $c_i=\lambda\min\big(\bar c, \frac{\pi(a_i|x_i)}{\mu(a_i|x_i)}\big)$. In the on-policy case, when $n=\infty$, V-trace then reduces to TD($\lambda$).
\end{remark}

\subsection{Actor-Critic algorithm}

\subsubsection*{Policy gradient}
In the on-policy case, the gradient of the value function $V^{\mu}(x_0)$ with respect to some parameter of the policy $\mu$ is
$$\nabla V^{\mu}(x_0) = \E_{\mu}\Big[ \textstyle{\sum_{s\geq 0}}\gamma^s \nabla\log\mu(a_s|x_s) Q^{\mu}(x_s, a_s)\Big],$$
where $Q^{\mu}(x_s, a_s)\eqdef \E_{\mu} \big[ \sum_{t\geq s}\gamma^{t-s} r_t |x_s,a_s \big]$ is the state-action value of policy $\mu$ at $(x_s,a_s)$. This is usually implemented by a stochastic gradient ascent that updates the policy parameters in the direction of 
$\E_{a_s\sim \mu(\cdot|x_s)}\Big[ \nabla\log\mu(a_s|x_s) q_s \big| x_s\Big]$, 
where $q_s$ is an estimate of $Q^{\mu}(x_s,a_s)$, and averaged over the set of states $x_s$ that are visited under some behaviour policy $\mu$. 


Now in the off-policy setting that we consider, we can use an IS weight between the policy being evaluated $\pi_{\bar\rho}$ and the behaviour policy $\mu$, to update our policy parameter in the direction of
\beq\label{eq:off-policy.PG}
\E_{a_s\sim \mu(\cdot|x_s)}\Big[ \frac{\pi_{\bar \rho}(a_s|x_s)}{\mu(a_s|x_s)} \nabla\log\pi_{\bar \rho}(a_s|x_s) q_s \big| x_s\Big]
\eeq
where $q_s \eqdef r_s + \gamma v_{s+1}$ is an estimate of  $Q^{\pi_{\bar \rho}}(x_s, a_s)$ built from the V-trace estimate $v_{s+1}$ at the next state $x_{s+1}$.
The reason why we use $q_s$ instead of $v_s$ as the target for our Q-value $Q^{\pi_{\bar \rho}}(x_s,a_s)$ is that, assuming our value estimate is correct at all states, i.e.~$V=V^{\pi_{\bar \rho}}$, then we have $\E [q_s|x_s,a_s] = Q^{\pi_{\bar \rho}}(x_{s},a_s)$ (whereas we do not have this property if we choose $q_t = v_t$). See Appendix~\ref{sec:analysis}
for analysis and Appendix~\ref{sec:appendix_q_estimation}
for a comparison of different ways to estimate $q_s$. 


In order to reduce the variance of the policy gradient estimate~\eqref{eq:off-policy.PG}, we usually subtract from $q_s$ a state-dependent baseline, such as the current value approximation $V(x_s)$.

Finally notice that \eqref{eq:off-policy.PG} estimates the policy gradient for $\pi_{\bar \rho}$ which is the policy evaluated by the V-trace algorithm when using a truncation level $\bar \rho$. However assuming the bias $V^{\pi_{\bar\rho}}-V^{\pi}$ is small (e.g.~if $\bar\rho$ is large enough) then we can expect $q_s$ to provide us with a good estimate of $Q^{\pi}(x_s,a_s)$. Taking into account these remarks, we derive the following canonical V-trace actor-critic algorithm.


\subsubsection*{V-trace actor-critic algorithm}
Consider a parametric representation $V_{\theta}$ of the value function and the current policy $\pi_{\omega}$. Trajectories have been generated by actors following some behaviour policy $\mu$. The V-trace targets $v_s$ are defined by \eqref{eq:target.off}. At training time $s$, the value parameters $\theta$ are updated by gradient descent on the $l2$ loss to the target $v_s$, i.e., in the direction of
$$\big( v_s - V_\theta(x_s)\big) \nabla_\theta V_\theta(x_s),$$
and the policy parameters $\omega$ in the direction of the policy gradient:
$$ \rho_s \nabla_\omega\log\pi_\omega(a_s|x_s) \big( r_s+\gamma v_{s+1} - V_\theta(x_s)\big).
$$
In order to prevent premature convergence we may add an entropy bonus, like in A3C, along the direction
$$-\nabla_\omega \sum_a \pi_\omega(a|x_s) \log\pi_\omega(a|x_s).$$
The overall update is obtained by summing these three gradients rescaled by appropriate coefficients, which are hyperparameters of the algorithm.







